\documentclass[conference]{IEEEtran}
\IEEEoverridecommandlockouts

\usepackage{cite}
\usepackage{amsmath,amssymb,amsfonts}
\usepackage{algorithmic}
\usepackage{graphicx}
\usepackage{textcomp}
\usepackage{xcolor}
\usepackage{booktabs}
\usepackage{multirow}
\usepackage{url}
\usepackage{hyperref}

\hypersetup{
    colorlinks=true,
    linkcolor=blue,
    citecolor=blue,
    urlcolor=blue
}

\def\BibTeX{{\rm B\kern-.05em{\sc i\kern-.025em b}\kern-.08em
    T\kern-.1667em\lower.7ex\hbox{E}\kern-.125emX}}

\begin{document}

\title{Investigating the Domain Generalization Gap in Neural Vulnerability Detection: From Human-Authored to LLM-Generated Source Code}

\author{
\IEEEauthorblockN{Anonymous Authors}
\IEEEauthorblockA{\textit{Department of Computer Science and Engineering} \\
\textit{Laboratory for Automated Program Analysis \& Software Vulnerability Research} \\
Artifact Repository: \url{https://github.com/0Tusharbhardwaj/nlp}}
}

\maketitle

\begin{abstract}
Automated source code synthesis driven by Large Language Models (LLMs) is transforming modern software engineering. However, recent empirical audits reveal that state-of-the-art code generators frequently introduce critical Common Weakness Enumerations (CWEs), such as buffer overflows (CWE-120), command injection (CWE-78), and memory management errors (CWE-415). While machine learning classifiers exhibit promising results in detecting vulnerabilities within human-authored codebases, whether these representations generalize to LLM-generated code remains largely unexamined.

In this paper, we conduct a systematic empirical investigation into the \textit{Domain Generalization Gap} of neural vulnerability detection. Differentiating our work from broad claims of ``whole-system security,'' we strictly focus on static, function-level binary vulnerability classification across three benchmarks comprising over 51,000 C functions: CodeXGLUE/Devign (human-authored commits from QEMU and FFmpeg), LLMVul (frontier LLM-generated code across 226 production GitHub repositories), and SeCodePLT (a curated benchmark of paired CVE vulnerabilities and ground-truth patches).

To eliminate synthetic evaluation artifacts, we enforce rigorous data hygiene: auditing repository leakage, constructing a strictly repository-disjoint holdout set ($0\%$ repository overlap), quarantining 2,679 ambiguous labels, and stripping comments via a syntax-preserving 5-state lexer. We evaluate classical lexical models (TF-IDF with Logistic Regression and Platt-calibrated Linear Support Vector Machines) against a contextual transformer (Microsoft CodeBERT, 125M parameters) fine-tuned on 21,854 Devign functions.

Our empirical findings demonstrate: (1) \textbf{In-Distribution Efficacy:} CodeBERT achieves 66.36\% accuracy, 64.71\% Macro F1, and 0.7301 ROC-AUC on human code, outperforming classical baselines; (2) \textbf{Domain Generalization Collapse:} When evaluated zero-shot on LLM-generated code ($N=4,120$), CodeBERT suffers a severe performance drop of $-56.48\%$ in Macro F1 (falling to 8.23\%) due to an acute false-positive surge triggered by stylistic distribution shift; (3) \textbf{Patch Discrimination Blindness:} On 2,480 balanced pairs of CVE vulnerabilities and patches (SeCodePLT), zero-shot models perform at near random guessing (50.00\% accuracy, 0.4993 ROC-AUC), proving that sequence-level pooled classifiers fail to capture localized, surgical vulnerability repairs; and (4) \textbf{Benchmark Memorization:} Devign's official test split shares 100\% of its repositories with its training set, indicating that reported in-distribution metrics partially reflect repository memorization. We release our complete reproducible artifact, model checkpoints, and evaluation suite.
\end{abstract}

\begin{IEEEkeywords}
Vulnerability Detection, Neural Code Models, CodeBERT, Domain Generalization, Large Language Models, Software Defect Prediction.
\end{IEEEkeywords}

\section{Introduction}

\subsection{The Proliferation of AI-Assisted Code}
Generative Large Language Models (LLMs)---such as OpenAI GPT-4, Anthropic Claude, and specialized code synthesis engines---are increasingly integrated into developer environments. Despite notable developer productivity gains, empirical evaluations have demonstrated the \textit{plausible correctness fallacy}: generated source code frequently satisfies superficial syntactic and functional requirements while embedding insidious Common Weakness Enumerations (CWEs)~\cite{pearce2022asleep}.

Because human review of large-scale automated code generation is cost-prohibitive, automated vulnerability detection is essential. Over the past decade, machine learning approaches have gained prominence, learning implicit statistical representations from code syntax and semantics~\cite{zhou2019devign, feng2020codebert, lu2021codexglue}. However, existing neural vulnerability detectors have predominantly been developed and validated on human-authored open-source software (e.g., historical commits in C/C++ projects). This assumes that representations learned on human code transfer effectively to synthetic, LLM-generated software---an assumption that lacks systematic empirical validation.

\subsection{Boundary of Scope: Vulnerability Detection vs. Security}
To establish methodological rigor, we explicitly distinguish \textit{static source code vulnerability detection} from \textit{software security}:
\begin{itemize}
    \item \textbf{Comprehensive Software Security} encompasses multi-tier defense mechanisms, including cryptographic key lifecycle management, identity and access control, dynamic runtime exploit mitigations (ASLR, stack canaries), network perimeter boundaries, and formal verification of distributed system states.
    \item \textbf{Our Focused Scope} is strictly confined to \textbf{static, function-level binary vulnerability classification}: determining whether an isolated C/C++ function exhibits structural or semantic patterns corresponding to known Common Weakness Enumerations (e.g., unchecked buffer copies, improper pointer deallocations).
\end{itemize}
We do not claim to provide a full-system security guarantee; rather, our objective is to measure the statistical reliability, calibration, and out-of-distribution transferability of neural code representations.

\subsection{Research Questions}
We guide our investigation through three central research questions:
\begin{itemize}
    \item \textbf{RQ1 (In-Distribution Performance):} How do classical lexical baselines (TF-IDF with Logistic Regression and Linear SVM) compare against deep contextual transformers (CodeBERT) on established human-authored C vulnerability benchmarks?
    \item \textbf{RQ2 (Cross-Domain Generalization):} Does a neural vulnerability classifier trained exclusively on human-written source code generalize to LLM-generated code under strict zero-leakage, repository-disjoint holdout conditions?
    \item \textbf{RQ3 (Surgical Patch Sensitivity):} Can sequence-level pooled transformer classifiers distinguish between vulnerable functions and their corresponding security patches in curated CVE benchmarks (SeCodePLT)?
\end{itemize}

\section{Related Work}

\subsection{Classical vs. Machine Learning Vulnerability Analysis}
Traditional Static Application Security Testing (SAST) tools (such as Flawfinder and Cppcheck) rely on deterministic Abstract Syntax Tree (AST) pattern matching and inter-procedural taint propagation. While sound within narrow rulesets, rule-based SAST exhibits high false-positive rates and cannot generalize to nuanced semantic variations.

Machine learning reformulates vulnerability identification as statistical sequence or graph classification. Early systems utilized bag-of-words and lexical token n-grams coupled with linear models or decision trees~\cite{scikit-learn}. Subsequent advances introduced Graph Neural Networks (e.g., Devign~\cite{zhou2019devign}), which extract composite code property graphs incorporating AST, control-flow, and data-flow representations.

\subsection{Pretrained Contextual Transformers for Code}
Transformer-based foundation models for code, notably Microsoft CodeBERT~\cite{feng2020codebert}, adapt the RoBERTa architecture~\cite{liu2019roberta} to programming languages. Pretrained on bimodal corpora spanning natural language documentation and six programming languages via Masked Language Modeling (MLM) and Replaced Token Detection (RTD), CodeBERT yields rich contextual representations. Fine-tuning CodeBERT via a sequence classification head on top of the leading classification token ($\mathbf{h}_{\text{<s>}}$) represents the benchmark standard across the CodeXGLUE defect detection suite~\cite{lu2021codexglue}.

\subsection{Vulnerability Characteristics of Generative Code}
Recent studies have highlighted the security vulnerabilities latent in generative code assistants~\cite{pearce2022asleep, secodeplt2024}. Generative models reproduce insecure idioms pervasive in public open-source training corpuses. Concurrently, AI-authored code displays distinctive syntactic regularities: uniform identifier conventions, standardized formatting, and repetitive structural boilerplate. Whether these stylistic idiosyncrasies confound neural classifiers trained on human idiosyncrasies remains an open empirical question.

\section{Dataset Hygiene, Provenance \& Partitioning}

Empirical machine learning research in software engineering requires rigorous data hygiene to prevent synthetic evaluation artifacts and repository memorization.

\begin{table*}[t]
\centering
\caption{Empirical Benchmark Dataset Inventory, Provenance, and Hygiene Partitions}
\label{tab:datasets}
\begin{tabular}{@{}lllrrp{5.5cm}@{}}
\toprule
\textbf{Benchmark} & \textbf{Language} & \textbf{Provenance} & \textbf{Samples} & \textbf{Vuln. Rate} & \textbf{Role in Empirical Evaluation} \\
\midrule
\textbf{CodeXGLUE / Devign} & C / C++ & QEMU, FFmpeg Commits & 27,318 & 45.8\% & Human training ($N=21,854$), validation ($2,732$), and test ($2,732$). \\
\textbf{LLMVul (Full)} & C / C++ & Production AI Code (226 Repos) & 21,430 & 8.8\% & Raw cross-domain AI-generated corpus. \\
\textbf{LLMVul (Clean Train)} & C / C++ & Greedy Disjoint (216 Repos) & 14,631 & 8.0\% & AI adaptation partition for cross-domain investigations. \\
\textbf{LLMVul (Holdout)} & C / C++ & Greedy Disjoint (8 Repos) & 4,120 & 8.8\% & Zero-leakage out-of-distribution evaluation cohort (E2). \\
\textbf{LLMVul (Gold Standard)} & C / C++ & Unanimous Human Expert Audit & 85 & 97.6\% & High-confidence ground-truth validation set. \\
\textbf{SeCodePLT} & C & Verified CVE Benchmarks & 2,480 & 50.0\% & 1,240 balanced pairs of vulnerable functions and patches. \\
\bottomrule
\end{tabular}
\end{table*}

\subsection{Benchmark Dataset Inventory}
As summarized in Table~\ref{tab:datasets}, we compile three diverse corpora comprising 51,228 functions:
\begin{enumerate}
    \item \textbf{CodeXGLUE / Devign:} 27,318 human-written C/C++ functions collected from QEMU and FFmpeg repositories with official train (21,854), validation (2,732), and test (2,732) partitions~\cite{zhou2019devign, lu2021codexglue}.
    \item \textbf{LLMVul:} 21,430 functions authored by frontier LLMs (GPT-4, Claude 3.5, GitHub Copilot) spanning 226 production GitHub repositories.
    \item \textbf{SeCodePLT:} NeurIPS 2024 benchmark providing 1,240 verified C Common Vulnerabilities and Exposures (CVEs) paired with their ground-truth patches ($N=2,480$)~\cite{secodeplt2024}.
\end{enumerate}

\subsection{Data Hygiene Audits}

\subsubsection{Internal Repository Overlap in Devign}
Our metadata audit of the standard Devign split revealed that the official test set shares \textbf{100\% of its source repositories} (`qemu` and `ffmpeg`) with the training set. Furthermore, exact text-hash deduplication identified:
\begin{itemize}
    \item 10 character-identical functions present in both training and test partitions.
    \item 57 functions in the training partition with conflicting labels (exact identical function text labeled as both vulnerable and safe).
\end{itemize}
This indicates that models evaluated on Devign may achieve inflated scores by memorizing repository-specific naming schemes and programming idioms rather than abstract vulnerability properties.

\subsubsection{Greedy Zero-Leakage Repository Partitioning (LLMVul)}
The raw LLMVul dataset contains 2,679 samples with ambiguous labels (`NaN`) caused by discordance among underlying automated verification scanners. We quarantined these samples into `llmvul_uncertain.jsonl` and excluded them from primary evaluation.

To enforce strict out-of-distribution evaluation, we developed a deterministic, greedy repository partitioner:
\begin{equation}
\mathcal{R}_{\text{train}} \cap \mathcal{R}_{\text{holdout}} = \emptyset
\end{equation}
The algorithm allocated 14,631 clean samples across 216 repositories to training and 4,120 clean samples across 8 complete, isolated repositories to the holdout cohort, ensuring \textbf{0\% repository leakage}.

\subsubsection{Syntax-Preserving Lexical Normalization}
Standard naive regular expression comment strippers (e.g., `re.sub(r'//.*', '', code)`) routinely corrupt string literals containing slash characters (e.g., URI schemes `"https://..."`). To prevent syntax distortion, we engineered a deterministic 5-state lexer:
\begin{equation}
S \in \{\text{NORMAL}, \text{IN\_STRING}, \text{IN\_CHAR}, \text{IN\_LINE\_CMT}, \text{IN\_BLOCK\_CMT}\}
\end{equation}
This state machine removes comments while preserving string literals, character constants, pointer dereferences, and macro boundaries.

\subsubsection{Token Truncation Audit}
Using the BPE tokenizer of `microsoft/codebert-base`, we audited the length distribution across all 21,854 Devign training functions. While the median length is 71 tokens, \textbf{23.12\% of functions (5,053 samples) exceed the maximum sequence limit of 512 tokens}. In our pipeline, we enforce a consistent head-truncation policy, retaining function declarations, formal parameter types, and initial control blocks while truncating trailing statements.

\section{Modeling Methodology \& Training Protocol}

\subsection{Classical Lexical Baselines}
We implement two classical baselines fitted on the 21,854 Devign training samples:
\begin{enumerate}
    \item \textbf{TF-IDF + Balanced Logistic Regression:} We extract word and syntax n-grams ($n \in [1, 2]$) up to 10,000 features. We optimize $L_2$-regularized cross-entropy with inverse regularization parameter $C=1.0$ via L-BFGS. Class weights are inversely balanced:
    \begin{equation}
    w_c = \frac{N}{2 \cdot N_c}
    \end{equation}
    where $N$ is total samples and $N_c$ is class frequency.
    \item \textbf{TF-IDF + Calibrated Linear SVM:} A Linear Support Vector Classifier (`LinearSVC`) optimizing squared-hinge loss with $L_2$ penalty ($C=1.0$). To produce calibrated posterior probabilities, we apply cross-validated \textbf{Platt scaling}~\cite{platt1999probabilistic}:
    \begin{equation}
    P(y=1 \mid x) = \frac{1}{1 + \exp\left(A \cdot f(x) + B\right)}
    \end{equation}
    where $f(x) = \mathbf{w}^T x + b$ is the uncalibrated decision margin.
\end{enumerate}

\subsection{Microsoft CodeBERT Fine-Tuning}
We fine-tune `microsoft/codebert-base` (12 layers, 768 hidden dimensions, 12 attention heads, 125M parameters). The classification head pools the leading classification token $\mathbf{h}_{\text{<s>}}$:
\begin{equation}
\mathbf{h}_{\text{pool}} = \tanh(\mathbf{W}_1 \mathbf{h}_{\text{<s>}} + \mathbf{b}_1)
\end{equation}
\begin{equation}
\hat{\mathbf{y}} = \text{softmax}(\mathbf{W}_2 \, \text{Dropout}(\mathbf{h}_{\text{pool}}) + \mathbf{b}_2)
\end{equation}

\subsubsection{Hyperparameters and Optimization}
Fine-tuning was executed on a Google Colab Tesla T4 GPU (16 GB VRAM) using Automatic Mixed Precision (FP16 via `torch.cuda.amp`). We employed the AdamW optimizer with weight decay $\lambda = 0.01$ and a peak learning rate of $\alpha = 2 \times 10^{-5}$ governed by a cosine learning rate scheduler. To prevent early gradient shock, we configured a linear warmup over the initial 10\% of iterations (341 warmup steps). Training proceeded for 3 complete epochs (4,098 total steps, batch size 16). Checkpoints were evaluated every 500 steps, and the model state exhibiting peak validation Macro F1 was preserved.

\section{Empirical Results \& Discussion}

\begin{table*}[t]
\centering
\caption{Empirical Vulnerability Classification Performance Across Experimental Cohorts}
\label{tab:main_results}
\begin{tabular}{@{}lllcccccc@{}}
\toprule
\textbf{Setting} & \textbf{Model} & \textbf{Evaluation Cohort ($N$)} & \textbf{Accuracy} & \textbf{Precision} & \textbf{Recall} & \textbf{Macro F1} & \textbf{FNR} & \textbf{ROC-AUC} \\
\midrule
\textbf{Baseline 1} & TF-IDF + Logistic Regression & Devign Test ($2,732$) & 61.35\% & 57.23\% & 62.79\% & 61.30\% & 37.21\% & 0.6734 \\
\textbf{Baseline 2} & TF-IDF + Linear SVM (Platt) & Devign Test ($2,732$) & 59.44\% & 56.93\% & 48.13\% & 58.48\% & 51.87\% & 0.6634 \\
\textbf{E1 (In-Distribution)} & \textbf{Microsoft CodeBERT} & \textbf{Devign Test ($2,732$)} & \textbf{66.36\%} & \textbf{68.96\%} & \textbf{48.69\%} & \textbf{64.71\%} & \textbf{51.31\%} & \textbf{0.7301} \\
\midrule
\textbf{E2 (AI Transfer)} & \textbf{Microsoft CodeBERT} & \textbf{LLMVul Holdout ($4,120$)} & \textbf{8.93\%} & \textbf{8.84\%} & \textbf{100.0\%} & \textbf{8.23\%} & \textbf{0.00\%} & \textbf{0.6888} \\
\textbf{Gold Standard} & Microsoft CodeBERT & Expert Consensus ($85$) & 97.65\% & 97.65\% & 100.0\% & 49.40\% & 0.00\% & 0.3675 \\
\textbf{CVE Patch Transfer} & Microsoft CodeBERT & SeCodePLT Pairs ($2,480$) & 50.00\% & 50.00\% & 99.84\% & 33.48\% & 0.16\% & 0.4993 \\
\bottomrule
\end{tabular}
\end{table*}

\begin{table}[htbp]
\centering
\caption{Confusion Matrix Breakdown Across Primary Experimental Settings}
\label{tab:cm_breakdown}
\begin{tabular}{@{}lrrrr@{}}
\toprule
\textbf{Evaluation Setting} & \textbf{TN} & \textbf{FP} & \textbf{FN} & \textbf{TP} \\
\midrule
\textbf{E1 (Devign In-Dist, $N=2,732$)} & 1,202 & 275 & 644 & 611 \\
\textbf{E2 (LLMVul AI Holdout, $N=4,120$)} & 4 & 3,752 & 0 & 364 \\
\textbf{SeCodePLT (CVE Pairs, $N=2,480$)} & 2 & 1,238 & 2 & 1,238 \\
\bottomrule
\end{tabular}
\end{table}

\subsection{RQ1: Human In-Distribution Performance (E1)}
As detailed in Table~\ref{tab:main_results}, CodeBERT outperforms classical lexical baselines on human-authored source code:
\begin{itemize}
    \item \textbf{Discriminative Power:} CodeBERT attains \textbf{66.36\% Accuracy} and \textbf{64.71\% Macro F1}, exceeding Logistic Regression (61.30\% F1, $+3.41\%$) and Linear SVM (58.48\% F1, $+6.23\%$).
    \item \textbf{ROC-AUC:} CodeBERT achieves \textbf{0.7301}, compared to 0.6734 for Logistic Regression and 0.6634 for Linear SVM ($+6.67\%$).
    \item \textbf{Precision and Noise Suppression:} CodeBERT yields a precision of \textbf{68.96\%} (275 false positives out of 1,477 non-vulnerable functions, Table~\ref{tab:cm_breakdown}), reflecting the capacity of multi-head self-attention to disambiguate benign buffer operations from unbounded copies.
\end{itemize}

\subsection{RQ2: Cross-Domain Generalization Collapse (E2)}
When CodeBERT is deployed zero-shot on LLM-generated code across the strictly disjoint holdout set ($N=4,120$), we observe a catastrophic collapse in discriminative ability:
\begin{equation}
\Delta \text{Macro F1} = -56.48\% \quad (64.71\% \rightarrow 8.23\%)
\end{equation}
\begin{equation}
\Delta \text{Accuracy} = -57.43\% \quad (66.36\% \rightarrow 8.93\%)
\end{equation}
As shown in Table~\ref{tab:cm_breakdown}, the model classified 3,752 out of 3,756 non-vulnerable AI functions as vulnerable, yielding only 4 true negatives. While Recall reaches 100\% (False Negative Rate = 0.0\%), this reflects a pathological loss of calibration: the model systematically flags AI-authored code as anomalous and vulnerable regardless of semantic safety. We attribute this to stylistic distribution shift: LLMs produce code with standardized naming, uniform indentation, and synthetic structure that human-trained sequence classifiers associate with aberrant code patterns.

\subsection{RQ3: Sensitivity to Surgical Security Patches (SeCodePLT)}
On the SeCodePLT benchmark comprising 2,480 balanced samples (1,240 verified CVE-vulnerable C functions paired with their ground-truth commit patches):
\begin{itemize}
    \item CodeBERT achieves exactly \textbf{50.00\% Accuracy} and an ROC-AUC of \textbf{0.4993}, performance that is statistically indistinguishable from a random coin toss.
    \item In 99.8\% of cases, the model predicted both the original vulnerable function and the patched function as vulnerable (1,238 false positives and 1,238 true positives, Table~\ref{tab:cm_breakdown}).
\end{itemize}
\textit{Analytical Takeaway:} Security patches in C typically involve minimal token edits---such as inserting an explicit bounds guard (`if (len >= MAX) return -1;`) or reordering a memory deallocation. Because sequence-level classification heads pool representations across all 512 input tokens, the 95\%+ unchanged token context dominates the sequence vector, rendering the localized patch invisible to the classifier.

\section{Threats to Validity}

\begin{enumerate}
    \item \textbf{Construct Validity (Function vs. Whole-Program Context):} Static analysis at function granularity ignores inter-procedural data flows, global state invariants, and execution wrappers. A function flagged as vulnerable in isolation may be safe due to input sanitization executed upstream.
    \item \textbf{Internal Validity (Sequence Truncation):} 23.12\% of functions exceed the 512-token boundary. Truncation discards trailing code, potentially removing the location of the actual vulnerability or its sanitizing check.
    \item \textbf{External Validity (Language Scope):} Our empirical results are established on C and C++ functions where memory safety (CWE-120, CWE-415) is paramount. In managed languages (e.g., Python, Go), cross-domain transfer dynamics may present different characteristics.
    \item \textbf{Benchmark Label Noise:} As documented in Section III, public vulnerability datasets contain labeling noise and contradictory annotations. We mitigated this by quarantining ambiguous samples and corroborating findings on an expert-audited gold cohort.
\end{enumerate}

\section{Artifact Availability \& System Architecture}

To ensure full reproducibility, we release our complete software pipeline as an open-source artifact at \url{https://github.com/0Tusharbhardwaj/nlp}:
\begin{itemize}
    \item \textbf{Trained Weights:} Full-scale baseline models and CodeBERT safetensors ($498$ MB) available via programmatic loaders.
    \item \textbf{Inference Service:} High-throughput FastAPI REST engine providing sub-15ms baseline inference and multi-model dispatch with transparent fallback attribution.
    \item \textbf{Interactive Dashboard:} React 19 and Monaco Editor interface offering preset CWE test templates and live confidence visualization.
    \item \textbf{Reproduction Suite:} End-to-end Python test suite (44 unit/integration tests) and self-contained Colab GPU notebooks.
\end{itemize}

\section{Conclusion \& Recommendations}

In this paper, we conducted a systematic empirical evaluation of neural vulnerability detection across human-authored and LLM-generated source code:
\begin{enumerate}
    \item Pretrained contextual transformers (CodeBERT) demonstrate superior in-distribution performance on human code (64.71\% Macro F1, 0.7301 ROC-AUC) over classical lexical baselines.
    \item Models trained exclusively on human-written software suffer catastrophic domain shift when evaluated zero-shot on AI-generated code ($-56.48\%$ Macro F1 drop), triggering a severe false-positive surge.
    \item Function-level pooled classifiers fail to detect subtle, surgical vulnerability patches, performing at random guessing levels on balanced CVE benchmarks (50.00\% accuracy on SeCodePLT).
    \item Standard benchmarks such as Devign suffer from significant internal repository leakage (100\% shared repositories between train and test splits).
\end{enumerate}
\textbf{Recommendation for Practice:} Software engineering teams deploying generative AI coding assistants should avoid relying on vulnerability detectors trained solely on historical human open-source code. Safe deployment necessitates domain-adapted fine-tuning on AI-synthesized corpora and hybrid coupling with dynamic sanitizers.

\bibliographystyle{IEEEtran}
\bibliography{references}

\end{document}

